\section*{Section 4 Concept Quiz}
\addcontentsline{toc}{section}{Section 4 Concept Quiz}

Use this as a short retrieval check after finishing the chapter. On the website, each item should accept an answer before revealing the hint and worked explanation.

\begin{bgcheck}{infinite-q1}{choice}{-infinity}[Concept check]
Evaluate \(\lim_{x\to2^-}\frac1{x-2}\).
\begin{bghint}
From the left, \(x-2\) is a small negative number.
\end{bghint}
\begin{bgreveal}
A positive numerator divided by a small negative number tends to \(-\infty\).
\end{bgreveal}
\end{bgcheck}

\begin{bgcheck}{vertical-q1}{integer}{3}[Concept check]
For \(f(x)=1/(x-3)^2\), enter the \(x\)-value of the vertical asymptote.
\begin{bghint}
A vertical asymptote occurs where the denominator is zero and does not cancel.
\end{bghint}
\begin{bgreveal}
\(x=3\).
\end{bgreveal}
\end{bgcheck}

\begin{bgcheck}{end-q1}{rational}{2/5}[Concept check]
Evaluate \(\lim_{x\to\infty}\frac{2x^2+1}{5x^2-3}\).
\begin{bghint}
The degrees match, so compare leading coefficients.
\end{bghint}
\begin{bgreveal}
The limit is \(2/5\).
\end{bgreveal}
\end{bgcheck}

\begin{bgcheck}{radical-infinity-q1}{integer}{-1}[Concept check]
Evaluate \(\lim_{x\to-\infty}\frac{\sqrt{x^2+1}}x\).
\begin{bghint}
Remember \(\sqrt{x^2}=|x|=-x\) when \(x<0\).
\end{bghint}
\begin{bgreveal}
The expression approaches \(-1\).
\end{bgreveal}
\end{bgcheck}

\begin{bgsummary}[title={Quiz use on the website}]
This page should report completion by concept, not merely a total score. A wrong answer should link back to the exact lesson that teaches the missing method. The same questions may appear in randomized numerical variants, but the canonical explanation should remain stable.
\end{bgsummary}
